\section{Related Work}\label{related-work}

\subsection{Probing for linguistic and semantic structure}\label{probing-for-linguistic-and-semantic-structure}

Linear probing \citep{alain2017probes} has become a standard tool
for reading off what information is encoded in neural network
representations. \citet{conneau2018probing} probed sentence
embeddings for syntactic properties; subsequent work probed for
part-of-speech, dependency relations, coreference, and world
knowledge. The linear representation hypothesis
\citep{park2024linear} formalizes the claim that concepts are
encoded as directions in representation space --- precisely the
assumption underlying our use of probe weight vectors as geometric
objects.

Two limitations of the probing literature motivate our approach.
First, most probing studies report only \emph{accuracy}, discarding the
learned probe parameters. We show that the probe weight vector ---
the normal to the classification hyperplane --- carries geometric
information about how concepts relate to each other. Second,
probing studies typically train one probe per concept, treating
concepts as independent. Our multi-probe geometric analysis
recovers inter-concept structure from independently trained probes.

\subsection{Geometry of concept representations}\label{geometry-of-concept-representations}

\citet{bolukbasi2016gender} demonstrated that gender bias in word
embeddings manifests as a geometric subspace, and that debiasing
amounts to projecting out a direction. \citet{arditi2024refusal}
showed that refusal behavior in instruction-tuned LLMs is mediated
by a single direction in activation space. These results establish
the precedent that high-level behavioral properties can be localized
to specific directions.

Our work extends this line from single directions to \emph{sets of
related directions}. Where prior work asked ``where is concept \(X\)?''
we ask ``what is the geometric relationship between concepts
\(X_1, \ldots, X_k\)?'' The cosine similarity matrix between
foundation probe directions is a form of representational similarity
analysis \citep[RSA;][]{kriegeskorte2008rsa} applied not to stimulus
response patterns but to the probes that decode them.

\subsection{Moral psychology and Moral Foundations Theory}\label{moral-psychology-and-moral-foundations-theory}

Moral Foundations Theory \citep[MFT;][]{haidt2012righteous,
graham2013mft} posits that human moral judgment draws on (at least)
five or six innate foundations: care/harm, fairness/cheating,
loyalty/betrayal, authority/subversion, sanctity/degradation, and
(later added) liberty/oppression. MFT predicts a structural
distinction between \emph{individualizing} foundations (care, fairness,
liberty), which protect individuals from harm, and \emph{binding}
foundations (loyalty, authority, sanctity), which bind individuals
into groups. This distinction has been validated in cross-cultural
surveys and predicts political orientation
\citep{graham2013mft}.

Alternative taxonomies exist. \citet{curry2019cooperation} propose
morality-as-cooperation, identifying seven moral domains grounded
in evolutionary game theory. Our use of MFT is pragmatic: the
six-foundation taxonomy provides a tractable set of directions for
geometric analysis, and the individualizing/binding prediction
provides a testable structural hypothesis.

\subsection{Moral reasoning in language models}\label{moral-reasoning-in-language-models}

\citet{hendrycks2021ethics} introduced the ETHICS benchmark for
evaluating moral reasoning in language models across justice,
deontology, virtue, utilitarianism, and commonsense dimensions.
Most work in this area evaluates models \emph{behaviorally} ---
measuring what outputs models produce in response to moral
scenarios. Our approach is \emph{representational}: we probe the
internal geometry of moral encoding, asking not whether the model
produces the right moral judgment but whether it has developed
structured representations of moral distinctions.

\subsection{Companion papers}\label{companion-papers}

This paper is the third in a series. Paper 1
\citep{reblitzrichardson2026fragility} established that moral
content is decodable by linear probes from the earliest layer
of OLMo-2, that probing accuracy saturates early during
pre-training, and that fragility testing --- injecting Gaussian
noise into activations --- resolves the continued development of
moral encoding after accuracy plateaus. Paper 2
\citep{reblitzrichardson2026dilution} extended the analysis to
MoE models (OLMoE-1B-7B), finding uniform moral encoding across
experts but a 74\(\times\) output scale gap that produces structural
fragility.

The present paper moves from \emph{binary} moral encoding (moral
vs.~neutral) to \emph{structured} moral encoding (framework-specific
directions and their inter-framework geometry). It reuses the
probing and fragility protocols from Papers 1 and 2, applying
them at the per-foundation level.
